

\subsection{Saint-Venant class}
The elastostatic field \((\bm u,\bm T)\) obeys \eqref{eq:3D-LE} with \(\bm T\,\bm n=\mathbf{0}\) on \(\BodyLateralSurface\).
We work in the relaxed setting: end tractions on \(\Section_0,\Section_L\) are not prescribed pointwise; only the resultants
\(\mathscr F(\bm u)\) and \(\mathscr M(\bm u)\) in \eqref{eq:def-resultants} are constrained by the data.

At this stage \emph{no} reduction is postulated for the in–plane block \(S_{\alpha\beta}\).
We nevertheless single out components that will drive the subsequent reductions:
\[
\AxialStress := \FrameBasis\cdot\bm T\,\FrameBasis \in \mathbb{R},
\qquad
\PointShearStress_\alpha := \FrameBasis_\alpha\cdot\bm T\,\FrameBasis \in \mathbb{R}\ (\alpha=1,2).
\]
Collect the shear components into the in–section vector \(\bm{\tau}:=\PointShearStress_\alpha\,\FrameBasis_\alpha\in\mathbb{R}^2\).
The Saint–Venant form
\begin{equation}
\bm T=\AxialStress\,\FrameBasis\otimes\FrameBasis+\FrameBasis\otimes\bm{\tau}+\bm{\tau}\otimes\FrameBasis
\label{eq:saint-venant-stress}
\end{equation}
is \emph{recovered} by the constructions in the two problems below; in particular \(S_{\alpha\beta}=0\) will follow from the field equations and boundary conditions under Ie\c{s}an’s postulate.

% --------------------------------------------------------------------

\subsection{Problem P1 extension--bending--torsion}

\subsubsection*{Problem statement}
Given resultants \((F_3,M_1,M_2,M_3)\) with \(F_\alpha=0\), determine \((\bm u,\bm T)\) such that:
\begin{enumerate}
  \item \textit{Partial continuum BVP.} In \(\Body\): \(\Div\,\StressTensor = \mathbf{0}\), \(\bm T=\lambda (\tr\StrainTensor) \mathbf{1} + 2G\,\StrainTensor\),
        \(\bm\varepsilon=\tfrac12(\nabla\bm u+\nabla\bm u^{\!\top})\); on \(\BodyLateralSurface\): \(\bm T\,\bm n=\mathbf{0}\).
        On \(\Section_0,\Section_L\) the traction fields are unknown pointwise, but satisfy the resultant constraints
        \[
        \FrameBasis\cdot\mathscr F(\bm u)=F_3,\qquad \mathscr{M}(\bm u)=\bm M,\qquad \FrameBasis_\alpha\cdot\mathscr F(\bm u)=0 .
        \]
  \item \textit{Ie\c{s}an kinematic postulate.} The axial derivative is rigid on each section:
        \begin{equation}
        \bm u'=\bm\alpha+\bm\beta\times\bm{r}
        \quad\Longleftrightarrow\quad
        \bm u'= a_{\varepsilon}\,\FrameBasis + \big(a_{\vartheta}\,\FrameBasis-\FrameBasis\times\bm a_\Section\big)\times\bm{r} ,
        \label{eq:Iesan-postulate-1}
        \end{equation}
        with constants \(\bm a_\Section\in\mathbb{R}^2\), \(a_\varepsilon,a_\vartheta\in\mathbb{R}\).
\end{enumerate}
This postulate fixes the \(x\)-dependence of \(\bm u\) and reduces the equilibrium system to planar problems on \(\Section\), from which \(S_{\alpha\beta}=0\) and \eqref{eq:saint-venant-stress} follow within the recovered solution.

\subsubsection*{Construction summary}
The torsion warping \(\WarpTwistIesan\) is the Neumann solution of
\begin{equation}
\Delta \WarpTwistIesan=0\ \text{ in }\Section,
\qquad
\partial_{\bm\nu}\WarpTwistIesan=\varepsilon_{\alpha\beta} n_\alpha x_\beta\ \text{ on }\SectionBoundary.
\label{eq:bvp-warp-twist-iesan}
\end{equation}
The displacement generated by \(\bm a=(\bm a_\Section,a_\varepsilon,a_\vartheta)\) is
\begin{equation}
\hat{\bm u}\{\bm a\}
=\Big(-\tfrac12 x^2\,\mathbf{1}-\nu\,\operatorname{dev}\bm{r}\otimes\bm{r}+x\,\FrameBasis\otimes\bm{r}\Big)\bm a_\Section
+(x\,\FrameBasis-\nu\,\bm{r})\,a_\varepsilon
+\big(x\,\FrameBasis\times\bm{r}+\WarpTwistIesan\,\FrameBasis\big)a_\vartheta .
\label{eq:iesan-displacement-1}
\end{equation}
The strain and the axial/shear components are
\begin{equation}
\bm E^{(a)}=
\big(a_\varepsilon+\bm{r}\cdot\bm a_\Section\big)\,(\FrameBasis\otimes\FrameBasis-\nu\,\mathbf{1}_\Section)
+\tfrac12\big(\FrameBasis\otimes\bm\varepsilon_\vartheta+\bm\varepsilon_\vartheta\otimes\FrameBasis\big),
\quad
\bm\varepsilon_\vartheta:=a_\vartheta\big(\FrameBasis\times\bm{r}+\nabla\WarpTwistIesan\big),
\label{eq:iesan-strain-tensor-1}
\end{equation}
\begin{equation}
\bm{\tau}=G\big(\nabla \WarpTwistIesan-\FrameBasis\times\bm{r}\big)a_{\vartheta},
\qquad
\AxialStress=E\big(\bm{r}\cdot\bm a_\Section +a_{\varepsilon}\big).
\label{eq:iesan-stress-1}
\end{equation}
The parameters are fixed by the resultants:
\begin{equation}
E\big(\IntRoR\,\bm a_\Section+A\,a_{\varepsilon}\,\CentroidLocation\big)=\FrameBasis\times\bm M,\qquad
EA\big(\bm a_\Section\cdot\CentroidLocation+a_{\varepsilon}\big)=-N,\qquad
G\,\Jo\,a_{\vartheta}=-T,
\label{eq:2-15}
\end{equation}
with
\begin{equation}
\Jo=\int_{\Section}\Big(r^2 + (\FrameBasis\times \bm{r})\cdot\nabla \WarpTwistIesan\Big)\,\mathrm{d}A.
\label{eq:def-j-origin}
\end{equation}

% --------------------------------------------------------------------

\subsection{Problem P2 flexure}

\subsubsection*{Problem statement}
Given a transverse resultant \(\ShearForce\) at \(\Section_1\) with \(F_3=0\) and \(\bm M=\mathbf{0}\), determine \((\bm u,\bm T)\) such that:
\begin{enumerate}
  \item \textit{Partial continuum BVP.} In \(\Body\): \(\Div\,\bm T=\mathbf{0}\), \(\bm T=\lambda(\tr\bm\varepsilon)\mathbf{1}+2G\,\bm\varepsilon\);
        \(\bm T\,\bm n=\mathbf{0}\) on \(\BodyLateralSurface\).
        On \(\Section_0,\Section_L\) the exact tractions are left unspecified, constrained only by
        \(\mathscr F(\bm u)=\ShearForce\) and \(\mathscr M(\bm u)=\mathbf{0}\) in \eqref{eq:def-resultants}, with \(\ShearForce\cdot\FrameBasis=0\).
  \item \textit{Ie\c{s}an kinematic postulate.} The axial derivative is rigid on each section as in \eqref{eq:Iesan-postulate-1}.
\end{enumerate}
The reduction implies \(\bm u'\) solves \(P_1\) with \(\bm M=\FrameBasis\times\ShearForce\), which fixes the amplitudes below.

\subsubsection*{Construction summary}
The flexure parameters solve
\begin{equation}
E\big(\IntRoR\,\bm b_{\Section}+A\,b_{\varepsilon}\,\CentroidLocation\big)=-\ShearForce,\qquad
\bm b_{\Section}\cdot\CentroidLocation+b_{\varepsilon}=0,\qquad
b_{\vartheta}=0.
\label{eq:2-18}
\end{equation}
The flexural warping \(\WarpShearIesan\) is given by
\begin{equation}
\left\{
\begin{aligned}
\Delta \WarpShearIesan&=-2\big((\bm{r}-\CentroidLocation)\cdot\bm b_\Section\big) && \text{in }\Section,\\
\partial_{\bm{\nu}} \WarpShearIesan&=\nu\,\bm{\nu}\cdot\Big(\operatorname{dev}\bm{r}\otimes\bm{r}-\bm{r}\otimes\CentroidLocation\Big)\bm b_\Section
&& \text{on } \SectionBoundary,\\
\int_{\Section}\WarpShearIesan\,\mathrm{d}A&=0, &&
\end{aligned}
\right.
\label{eq:bvp-warp-shear-iesan}
\end{equation}
and the displacement is
\begin{equation}
\begin{aligned}
\hat{\bm{u}}_{\Section}(x,\bm{r})
&=-\tfrac16\,x^3\,\bm{b}_{\Section}
+x\,\nu\Big(\tfrac12 r^{2}\,\bm{b}_{\Section}-(\bm{b}_{\Section}\cdot\bm{r})\,\bm{r}\Big)
+x\,\nu\,(\CentroidLocation\cdot\bm b_\Section)\,\bm{r}
+c_{\vartheta}\, x\,(\FrameBasis\times \bm{r}),
\\
\FrameBasis\cdot\hat{\bm{u}}(x,\bm{r})
&=\tfrac12\,x^2\,\big((\bm{r}-\CentroidLocation)\cdot\bm b_{\Section}\big)
+c_{\vartheta}\,\WarpTwistIesan(\bm{r})+\WarpShearIesan(\bm{r}).
\end{aligned}
\label{eq:iesan-displacement-2}
\end{equation}
The scalar \(c_{\vartheta}\) is determined by
\begin{equation}
\Jo \,c_{\vartheta}
= -\int_{\Section} \Big[(\FrameBasis\times \bm{r})\cdot\nabla \WarpShearIesan
+ \tfrac{1}{2}\,\nu\, r^2\,(\FrameBasis\times \bm{r})\cdot\bm b_{\Section}\Big]\,\mathrm{d}A .
\label{eq:iesan-2-22}
\end{equation}
The strain and the Saint–Venant components are
\begin{equation}
\begin{aligned}
\bm E(\hat{\bm u})
&=(\varepsilon\,\FrameBasis+\bm{\varepsilon}_{\gamma})\stackrel{\rm s}{\otimes}\FrameBasis-\nu\,\varepsilon\,\mathbf{1}_\Section,\\
\varepsilon&=x\,(\bm{r}-\CentroidLocation)\cdot\bm b_\Section,\\
\bm{\varepsilon}_{\gamma}
&=(a_\vartheta+c_\vartheta)\big(\nabla \WarpTwistIesan + \FrameBasis\times \bm{r}\big)
-\nu\Big(\operatorname{dev}(\bm{r}\otimes\bm{r})-\bm{r}\otimes\CentroidLocation \Big)\bm b_{\Section}
+\nabla (\bm{\WarpShearIesan}\cdot\bm b_{\Section}),
\end{aligned}
\label{eq:iesan-strain-2}
\end{equation}
\begin{equation}
\AxialStress(\hat{\bm u})=E\,x\,(\bm{r}-\CentroidLocation)\cdot\bm b_\Section,
\qquad
\bm{\tau}(\hat{\bm u})
=G\Big[(a_{\vartheta}+c_{\vartheta})\big(\nabla \WarpTwistIesan + \FrameBasis\times \bm{r}\big)
-\nu\Big(\operatorname{dev}(\bm{r}\otimes\bm{r})-\bm{r}\otimes\CentroidLocation \Big)\bm b_{\Section}
+\nabla (\bm{\WarpShearIesan}\cdot\bm b_{\Section})\Big].
\label{eq:iesan-stress-2}
\end{equation}
In both \(P_1\) and \(P_2\) the recovered solution satisfies \(S_{\alpha\beta}=0\), i.e. it falls within \eqref{eq:saint-venant-stress}.

% --------------------------------------------------------------------

\subsection{Saint-Venant equilibrium identities}
For any field of the form \eqref{eq:saint-venant-stress} with \(\bm T\,\bm n=\mathbf{0}\) on \(\BodyLateralSurface\),
the sectional equilibrium reduces to
\begin{subequations}
\label{eq:sv-equilibrium}
\begin{align}
\bm{\tau}'&=\mathbf{0} && \text{in }\Section,\\
\DivS \bm{\tau}+\AxialStress'&=0 && \text{in }\Section,\\
\bm{\tau}\cdot\bm{\nu}&=0 && \text{on } \SectionBoundary.
\end{align}
\end{subequations}
Thus \(\bm{\tau}\) is \(x\)-independent, and for flexure
\(\DivS\bm{\tau}=-E\,(\bm{r}-\CentroidLocation)\cdot\bm b_\Section\).

% --------------------------------------------------------------------

\subsection{General Saint-Venant solution}
Combining \eqref{eq:iesan-displacement-1} and \eqref{eq:iesan-displacement-2} yields
\begin{equation}
\begin{aligned}
\hat{\bm u}\{\bm a,\bm b_\Section\}
&=\Big(-\tfrac12 x^2\,\mathbf{1}-\nu\,\operatorname{dev}\bm{r}\otimes\bm{r}+x\,\FrameBasis\otimes\bm{r}\Big)\bm a_\Section
+(x\,\FrameBasis-\nu\,\bm{r})\,a_\varepsilon
+\big(x\,\FrameBasis\times\bm{r}+\WarpTwistIesan\,\FrameBasis\big)a_\vartheta\\
&\quad-\tfrac16\,x^3\,\bm b_\Section
-x\,\nu\Big(\operatorname{dev}\bm{r}\otimes\bm{r}-\bm{r}\otimes\CentroidLocation\Big)\bm b_\Section
+\tfrac12 x^2\,\FrameBasis\otimes(\bm{r}-\CentroidLocation)\,\bm b_\Section
+\FrameBasis\otimes\bm{\WarpShearIesan}(\bm{r})\,\bm b_\Section\\
&\quad+\big(x\,\FrameBasis\times\bm{r}+\WarpTwistIesan\,\FrameBasis\big)c_\vartheta,
\end{aligned}
\label{eq:iesan-general-solution}
\end{equation}
with amplitudes fixed by \eqref{eq:2-15} and \eqref{eq:2-18}, \(c_\vartheta\) by \eqref{eq:iesan-2-22},
and warpings by \eqref{eq:bvp-warp-twist-iesan}–\eqref{eq:bvp-warp-shear-iesan}.
